\documentclass[12pt]{article}

% Pachete necesare pentru codificarea caracterelor și diacritice
\usepackage[utf8]{inputenc}
\usepackage[romanian]{babel}

% Pachete matematice avansate
\usepackage{amsmath}
\usepackage{amsfonts}
\usepackage{amssymb}
\usepackage{amsthm}

% Setarea marginilor paginii
\usepackage[margin=1in]{geometry}

\title{Seminar 7: Criptosistemul RSA -- Rezolvare Problemă}
\author{}
\date{} % Elimină data și spațiul nefolosit sub titlu

\begin{document}

\maketitle

\section*{Problemă: Criptosistemul RSA (Percy și Charlie)}

\textit{Percy și Charlie comunică folosind criptosistemul RSA. Percy are cheia publică: $n = 187$ și $e = 107$.}
\begin{enumerate}
    \item[a.] \textit{Aflați cheia privată a lui Percy.}
    \item[b.] \textit{Charlie îi transmite lui Percy mesajul \texttt{ABACFPFP}. Știind că lungimea blocurilor mesajelor în clar este 1 și a mesajelor criptate este 2, decriptați textul.}
\end{enumerate}



\subsection*{a. Determinarea Cheii Private a lui Percy}

Cheia privată a lui Percy este perechea $(n, d)$, unde $d$ este exponentul de decriptare. Pentru a-l calcula pe $d$, trebuie mai întâi să factorizăm modulul public $n = 187$.

\subsubsection*{Pasul 1: Factorizarea lui $n$}
Căutăm numerele prime $p$ și $q$ astfel încât $n = p \cdot q$. Observăm prin verificări directe că:
\[
187 = 11 \times 17 \implies p = 11, \quad q = 17
\]

\subsubsection*{Pasul 2: Calculul funcției indicatoare a lui Euler $\varphi(n)$}
\[
\varphi(n) = (p - 1)(q - 1) = (11 - 1)(17 - 1) = 10 \times 16 = 160
\]

\subsubsection*{Pasul 3: Calculul exponentului de decriptare $d$}
Exponentul $d$ este inversul modular al lui $e = 107$ modulo $\varphi(n) = 160$, adică:
\[
107 \cdot d \equiv 1 \pmod{160}
\]
Aplicăm algoritmul lui Euclid extins pentru perechea $(160, 107)$:
\begin{align*}
160 &= 107 \times 1 + 53 \\
107 &= 53 \times 2 + 1
\end{align*}

Reconstituim restul $1$ prin substituție inversă:
\begin{align*}
1 &= 107 - 53 \times 2 \\
1 &= 107 - (160 - 107 \times 1) \times 2 \\
1 &= 3 \times 107 - 2 \times 160
\end{align*}
Trecând relația modulo $160$, obținem $3 \times 107 \equiv 1 \pmod{160}$. Prin urmare, $d = 3$.

\textbf{Concluzie pct. a:} Cheia privată a lui Percy este $\mathbf{(n = 187, d = 3)}$.

---

\subsection*{b. Decriptarea Mesajului}

Mesajul cifrat primit este \texttt{ABACFPFP}. Știm că blocurile criptate au lungimea $2$ (digrafi), iar blocurile în clar au lungimea $1$ (monoradix).

\subsubsection*{Pasul 1: Segmentarea și conversia numerică în baza 26}
Împărțim textul cifrat în blocuri de câte 2 caractere și le transformăm în valori numerice în baza $26$, utilizând formula $c = \text{cod}(C_1) \cdot 26 + \text{cod}(C_2)$:
\begin{enumerate}
    \item \texttt{AB} $\rightarrow (0, 1)_{26} = 0 \cdot 26 + 1 = 1$
    \item \texttt{AC} $\rightarrow (0, 2)_{26} = 0 \cdot 26 + 2 = 2$
    \item \texttt{FP} $\rightarrow (5, 15)_{26} = 5 \cdot 26 + 15 = 130 + 15 = 145$
    \item \texttt{FP} $\rightarrow (5, 15)_{26} = 145$
\end{enumerate}

\subsubsection*{Pasul 2: Aplicarea formulei de decriptare RSA}
Pentru fiecare valoare criptată $c$, blocul în clar sub formă numerică $m$ se obține prin calculul:
\[
m \equiv c^d \pmod{n} \implies m \equiv c^3 \pmod{187}
\]

\begin{itemize}
    \item \textbf{Pentru primul bloc ($c_1 = 1$):}
    \[ m_1 \equiv 1^3 = 1 \pmod{187} \rightarrow \mathbf{1} \]
    
    \item \textbf{Pentru al doilea bloc ($c_2 = 2$):}
    \[ m_2 \equiv 2^3 = 8 \pmod{187} \rightarrow \mathbf{8} \]
    
    \item \textbf{Pentru al treilea bloc ($c_3 = 145$):}
    Calculăm $145^3 \pmod{187}$. Reducem mai întâi baza observând că:
    \[ 145 \equiv 145 - 187 = -42 \pmod{187} \]
    Atunci:
    \[ 145^2 \equiv (-42)^2 = 1764 \pmod{187} \]
    Efectuăm împărțirea: $1764 = 187 \times 9 + 81 \implies 145^2 \equiv 81 \pmod{187}$.
    
    Acum înmulțim cu încă un factor:
    \[ 145^3 \equiv 145^2 \cdot 145 \equiv 81 \cdot (-42) = -3402 \pmod{187} \]
    Efectuăm împărțirea valorii absolute: $3402 = 187 \times 18 + 36 \implies -3402 \equiv -36 \pmod{187}$.
    
    Aducem la rest pozitiv:
    \[ m_3 = -36 + 187 = \mathbf{151} \]
    
    \item \textbf{Pentru al patrulea bloc ($c_4 = 145$):}
    Fiind identic cu blocul anterior, obținem direct aceiași valoare:
    \[ m_4 = \mathbf{151} \]
\end{itemize}

\subsubsection*{Pasul 3: Transpunerea valorilor $m$ în caractere (baza 26, lungime 1)}
Deoarece lungimea blocurilor în clar este 1, fiecare valoare obținută reprezintă codul unei singure litere modulo 26 ($m \pmod{26}$):
\begin{enumerate}
    \item $m_1 = 1 \implies 1 \pmod{26} \rightarrow \mathbf{B}$
    \item $m_2 = 8 \implies 8 \pmod{26} \rightarrow \mathbf{I}$
    \item $m_3 = 151 \implies 151 = 26 \times 5 + 21 \implies 21 \pmod{26} \rightarrow \mathbf{V}$
    \item $m_4 = 151 \implies 21 \pmod{26} \rightarrow \mathbf{V}$
\end{enumerate}

\subsection*{Concluzie Finală}
Prin concatenarea caracterelor rezultate din decriptarea celor 4 blocuri succesive, textul în clar obținut este:
\[
\text{Mesajul decriptat: } \mathbf{\texttt{BIVV}}
\]

\end{document}